% Source: https://artofproblemsolving.com/community/c3h1600011_acute_and_obtuse_angles
% Stable identity: aops-topic-1600011-q01
Let each acute angle be $a$ and each obtuse angle be $b$ (vertical angles are equal). Adjacent angles are supplementary, so $a+b=180^\circ$, and the condition is $2a=\dfrac {b}{4}$, i.e.\ $b=8a$. Then $9a=180^\circ$, so
\[
a=20^\circ,\qquad b=160^\circ,
\]
and the four angles are $\boxed{20^\circ,\ 160^\circ,\ 20^\circ,\ 160^\circ}$. Check: $20^\circ+20^\circ=40^\circ=\dfrac{1}{4}\cdot160^\circ$.

\begin{asy}
usepackage("amsmath");
size(5cm);
defaultpen(fontsize(10pt));
real t=20;
pair O=(0,0);
draw(3*dir(180)--3*dir(0),linewidth(1));draw(3*dir(180+t)--3*dir(t),linewidth(1));
draw(arc(O,1.3,0,t),red);label("$20^\circ$",1.75*dir(t/2),red);
draw(arc(O,1.3,180,180+t),red);label("$20^\circ$",1.75*dir(180+t/2),red);
draw(arc(O,0.6,t,180),blue);label("$160^\circ$",0.95*dir(100),blue);
draw(arc(O,0.6,180+t,360),blue);label("$160^\circ$",0.95*dir(280),blue);
label("\mbox{}",(min(currentpicture,user=true)+max(currentpicture,user=true))/2,black); // keeps the text after the figure black when LaTeX typesets the labels inline
\end{asy}
